% RECOVERED_RESULT from II REV09. See MAPA_PROCEDENCIA.json.
% Selection of lines and explicit mechanical changes; no new proof.
\section{Action, gravity, and thermal information}
\label{v:ant:sec:gravity}
\subsection{Dodecaphasic period and circular realization of action}
The longitudinal composition in~\eqref{xii:radial:length} determines
\(L=L_*\alpha^{16}5759/23040\), and its compatible clock satisfies
\(t_0=\pi L/(54c_{\rm int})\). The calendar already constructed determines
$T_{108}=108t_0$. This section transports that realization to the
circular reader, retaining \(R_{108}=L\).
The constitutive velocity, duration, and path length are
obtained from the identities in
Proposition~\ref{iii:prop:identidad-velocidad} and
equations~\eqref{iii:eq:velocidad-longitud-reloj} and
\eqref{iii:eq:longitud-elemental-constitutiva}.
In its circular realization,
\[
\omega_{108}=\frac{2\pi}{108t_0},\qquad
R_{108}=\frac{c_{\rm int}}{\omega_{108}}
=\frac{54}{\pi}\ell_0,\qquad \ell_0=c_{\rm int}t_0.
\]
The unit $\ell_0$ and the time and velocity charts remain
explicit; the SI chart does not select the cycle.
For each oriented action section $a\in\{\mathrm{pre},\mathrm{ret}\}$,
\[
E_{108,a}=\hbar_a\omega_{108},\qquad
m_{108,a}=\frac{E_{108,a}}{c_{\rm int}^2}.
\]
The reduced and full lengths are, respectively,
\[
\frac{\hbar_a}{m_{108,a}c_{\rm int}}=R_{108},
\qquad
\frac{h_a}{m_{108,a}c_{\rm int}}=2\pi R_{108}.
\]
The same action and duration enter both readouts.

\begin{proposition}[Radius-coincidence selector]
In the dimensional collection
\[
G_a(\vartheta)=\vartheta\frac{c_{\rm int}^5t_0^2}{\hbar_a},
\qquad \vartheta>0,
\]
the circular-realization condition
$G_a(\vartheta)m_{108,a}/c_{\rm int}^2=R_{108}$
has the unique solution
\[
\vartheta^\circ_{108}=(54/\pi)^2.
\]
\end{proposition}
\begin{proof}
The first radius is
$r_g(\vartheta)=\vartheta(\pi/54)\ell_0$.
Comparing it with $(54/\pi)\ell_0$, with $\ell_0>0$, reduces the
equality to a strictly increasing linear equation.
\end{proof}
Circular coincidence is the restriction to the unit mode of the positive
radial reader in Theorem~\ref{xii:radial:rigidity}. Its solution
retains the longitudinal and metric compositions declared there.
The reduced radius is \(Gm/c^2\); the Schwarzschild radius is
\(2Gm/c^2\) and carries its corresponding factor of two.

\begin{proposition}[Reciprocal action and gravity, conserved area]
In the two preceding charts,
\[
\frac{G_{\rm pre}}{G_{\rm ret}}=R_{\rm act}^{-1},\qquad
\hbar_aG_a=\vartheta^\circ_{108}c_{\rm int}^5t_0^2,
\qquad
\ell_P^2=\frac{\hbar_aG_a}{c_{\rm int}^3}
=\vartheta^\circ_{108}\ell_0^2.
\]
\end{proposition}
\begin{proof}
The definition of $G_a$ contains $\hbar_a^{-1}$, while all remaining factors
are common to the two charts. Taking the ratio and product yields
the three equalities.
\end{proof}
This reciprocity explains why orientation information can
vary while the area remains unchanged. It thus provides a precise relation
between the reciprocal action sections, the Planck scale, and the area operator
of Section~\ref{v:ant:sec:area}.
